\begin{figure}[tbp]
\centering
\ViewHeading{(a) The first component already gives the signature metric}
\begin{tikzpicture}
\begin{axis}[viewaxis,width=12.5cm,height=4.6cm,xmin=.5,xmax=24.5,
 xlabel={$q$},ylabel={$q^{-1}|v_q(x)-v_q(y)|$},ymin=0,ymax=1.7,legend pos=north east]
\addplot[viewblue,ycomb,mark=*,mark size=1.4pt] table[x=q,y=component] {\ArithmeticSignatureData};
\addlegendentry{$x=0.17,\ y=0.43$}
\addplot[vieworange,dashed] table[x=q,y=first_component] {\ArithmeticSignatureData};
\addlegendentry{$|v_1(x)-v_1(y)|$}
\end{axis}
\end{tikzpicture}
\ViewHeading{(b) Scalar costs forget information retained by the signature}
\begin{tikzpicture}
\node[vbox] (positive) at (-4,0) {$x$};
\node[vbox] (negative) at (-4,-1.2) {$-x$};
\node[vinput,text width=4.7cm] (scalar) at (1,-.6) {same scalar profile\\$(b_q(x))_q=(b_q(-x))_q$};
\draw[varrow] (positive)--(scalar);
\draw[varrow] (negative)--(scalar);
\node[align=center,font=\footnotesize,text width=3.4cm] at (5.1,-.6)
 {If $x\neq-x$, then\\$U(x)\neq U(-x)$.};
\end{tikzpicture}
\caption{Panel (a) samples the general inequality proved in the text;
the equality of the supremum with the first component holds at every
scale, not just the 24 plotted ones. The scalar profile determines costs
but is not an injective parameterization. The full marked signature is.}
\label{fig:arithmetic-signature}
\end{figure}